% GY51310C
% Versorgung des Neugeborenen
% 14.0
% 2013-11-30

.. _m-neugeborenes:


\begin{itemize}
  -   Neugeborenes abrubbeln
  -   Erstbeurteilung:
  \begin{itemize}
    -   Muskeltonus
    -   Atemweg
    \begin{itemize}
      -   Bei Atemwegsverlegung: Mittels Oro-Sauger absaugen
      \hfill\Commentary[Absaugen Neugeborener]{ASBÖ
      Lehrmeinung 18/2011 \enquote{Absaugen Neugeborener im
      Rettungsdienst}:cite:`AsboeLehrmeinungRd2011-19`\ : \enquote{Mehrere
      Studien zeigten keine Verbesserung im Outcome
      gegenüber abgesaugter Neugeborener. Daher ist es
      obligat, keine Verzögerung bei den BLS-Maßnahmen
      hervorzurufen, welche durch lange Absaugversuche
      zustande kommen können. Lediglich die
      Atemwegsobstruktion durch Mekonium muss behandelt
      werden, indem die Atemwege freigemacht werden. Selbst
      missfarbiges Fruchtwasser stellt keine Indikation zur
      Absaugung
      dar.}}
      \ParBugfix
      ERC 2010:cite:`Erc2010Sek06De`\ : \enquote{Absaugen ist
      nur notwendig, wenn die Atemwege verlegt sind. Eine
      solche Verle- gung kann aufgrund von Mekonium, auch
      wenn das Neugeborene keine Mekoniumablagerungen auf
      der Haut zeigt, aber auch Blutkoageln, zähem Schleim
      oder Vernix bestehen. Wird abgesaugt, ist zu bedenken,
      dass zu heftiges oropharyngeales Ab- saugen das
      Einsetzen einer suffizienten Spontanatmung verzögern
      und zu einem Laryngospasmus sowie zu einer vagalen
      Bradykardie führen kann. Das Vorhandensein von
      zähflüssigem Mekonium beim schlaffen, avitalen
      Neugeborenen ist die einzige Situation, in der ein
      sofortiges Absaugen des Oropharynx zu erwägen ist.}}
      }%Commentary
    \end{itemize}
    -   Atmung
    
    Bei Schnappatmung oder fehlender Atmung: Atemwege öffnen
    5 initiale Beatmungen
    -   Kreislauf: \Abk{HF} messen (durch Auskultation
    an der Herzspitze (nur geübtes Personal) oder an
    *Oberarmarterie*). Wenn ein *geeignetes
    Pulsoxymeter* vorhanden ist, kann auch dieses verwendet werden.
    \hfill\Commentary[Ort der HF-Messung beim Neugeborenen]
    {ERC 2010:cite:`Erc2010Sek06De`\ : \enquote{Die beste Methode
    zur Beurteilung der Herzfrequenz ist die direkte Auskultation mit dem Stethoskop über der Herzspitze. Das Tasten des
    Pulses an der Basis der Nabelschnur ist oft möglich, kann
    aber irreführend sein. Eine Beurteilung der Herzfrequenz
    allein über die Pulsation der Nabelschnur ist nur
    zuverlässig, wenn die Herzfrequenz über 100\PerMin* liegt.}}
    \ParBugfix
    Bei Sanitätern, insbesondere der
    niedrigeren Ausbildungsstufen, kann nicht davon
    ausgegangen werden, dass eine Auskultation in dieser
    stressreichen Situation fehlerfrei und zuverlässig angewendet werden kann. Die
    Oberarmarterie erscheint in diesem Fall als
    pragmatischer Kompromiss.}
  \end{itemize}
\end{itemize}


%\TableWide
%[]%1 Float
%{}%2 Titel
%{}%3 Beschreibung
%{}%4 Label/Hook
%{}%5 Quelle
%{}%6 Lizenz
%{}%7 Tabelle
\TableWide
[]%1 Float
{Versorgung des Neugeborenen abhängig vom Befund}
{}
{\label{fig:neugeborenes-versorgung}}
{}
{}
{%
\begin{tabularx}{\linewidth}{XX}
\toprule\addlinespace
\multicolumn{1}{c}{\includegraphics[width=1cm]{./icons/sm-basic.png}}
&
\multicolumn{1}{c}{\includegraphics[width=1cm]{./icons/sm-pale.png}~~\includegraphics[width=1cm]{./icons/sm-pale-frown.png}}
\\\addlinespace
\noindent\TabHeading{Kräftiges Schreien/suffiziente Atmung}
&
\noindent\TabHeading{Insuffiziente Spontanatmung,
Atemstillstand oder Bradykardie}
\\\addlinespace\cmidrule(lr){1-1}\cmidrule(lr){2-2}\addlinespace
\begin{itemize}
  \item[\eye] Guter Muskeltonus,
  \item[\eye] Herzfrequenz \textgreater 100\PerMin*
\end{itemize}




&
\begin{itemize}
  \item[\eye] Normaler, reduzierter Muskeltonus, oder sogar
  schlaffer Muskeltonus (\enquote{floppy})
  \item[\eye] Herzfrequenz \textless 100\PerMin*; Bradykarde
  oder nichtnachweisbare Herzfrequenz,
  \item[\eye] Oft mit ausgeprägter Blässe als Zeichen
  einer schlechten Perfusion.
\end{itemize}

\\\addlinespace\cmidrule(lr){1-1}\cmidrule(lr){2-2}\addlinespace
\begin{itemize}
  -   Abtrocknen und Einwickeln in warme Tücher:Warm
  einpacken (nur das Gesicht darf frei bleiben) und der Mutter
  überreichen.
  -   Abnabeln, frühestens nach 1\Min*:
  \begin{enumerate}
    -   Erst wird eine Handbreite vom Kind
    entfernt die erste Nabelklemme gesetzt.
    -   Dann wird die Nabelschnur Richtung Mutter ausgestreift und
    wieder eine Handbreit entfernt die zweite Klemme gesetzt.
    -   Zwischen den beiden
    Klemmen wird die Nabelschnur durchtrennt.
  \end{enumerate}
  -   Das Neugeborene
  kann der Mutter auf den Bauch gelegt werden. Durch den
  Kontakt zur Haut der zugedeckten Mutter wird das Baby
  gewärmt.
  
  -   Es kann zu diesem Zeitpunkt bereits an die Brust
  angelegt werden.
\end{itemize}
  &
  % Nachdem ein Neugeborenes dieser Gruppe getrocknet und in
  % warme Tücher gewickelt wurde, ist meist eine kurze
  % Maskenbeatmung ausreichend. Kommt es darunter nicht zu einem
  % adäquaten Anstieg der Herzfrequenz, müssen Herzdruckmassagen
  % durchgeführt werden.
  
  \begin{itemize}
    %     -   O₂-Gabe (\enquote{O₂-Dusche}) für 30 Sek.
    -   Atemwege freimachen, absaugen
    -   5 Beatmungen
    -   Wenn keine Steigerung der Herzfrequenz:
    \begin{itemize}
      %     -   O₂-Gabe (\enquote{O₂-Dusche}) für 30 Sek.
      -   Lagerung und Atemwege überprüfen
      -   5 Beatmungen wiederholen
    \end{itemize}
    -   Wenn weiterhin keine ausreichende Atmung oder
    \EE{Herzfrequenz \textless 60\PerMin*}:
    \begin{itemize}
      -   \Sign{→} Reanimation des Neugeborenen
    \end{itemize}
    \FullPrettyRef{m:neugeborenes-basisreanimation}
  \end{itemize}
  
  % Nachdem ein Neugeborenes dieser Gruppe getrocknet und in
  % warme Tücher gewickelt wurde, müssen unverzüglich die
  % Atemwege geöffnet werden. Die Lungen müssen belüftet werden,
  % und das Kind muss beatmet werden. Wurden die Atemwege
  % erfolgreich geöffnet und die Lungen belüftet, können
  % außerdem Herzdruckmassagen und möglicherweise eine
  % Medikamentengabe notwendig sein. Eine sehr kleine Gruppe von
  % Neugeborenen bleibt trotz adäquater Spontanatmung und guter
  % Herzfrequenz hypox­ ämisch. Die Diagnosen umfassen hier ein
  % weites Spektrum, wie z. B. eine Zwerchfellhernie,
  % Surfactant-Mangel, eine kongenitale Pneumonie, einen
  % Pneumothorax oder ein angeborenes zyanotisches Herzvitium.
  
  
  \\\addlinespace\bottomrule
\end{tabularx}
}


:cite:`Erc2010Sek06De,AsboeLehrmeinungRd2011-19,AsboeLehrmeinungRd2011-18`\ 